\section{Condições implementadas}
Escreva $P=P(\alpha)$, $P_+=P(\alpha_+)$ e $\eta=\gamma^2$.
No caso congelado, $P_+=P$ e $\dot P=0$.
Minimiza-se $\eta$ sujeito a $P\succ0$ e às condições abaixo, impostas por
positividade dos coeficientes polinomiais. Trata-se de uma condição suficiente,
que pode ser conservadora. Não se afirma obter a norma robusta exata.

Para $H_\infty$ contínuo:
\[
\begin{bmatrix}
-A^TP-PA-\dot P&-PB&-C^T\\
-B^TP&\eta I&-D^T\\-C&-D&I
\end{bmatrix}\succ0.
\]
Para $H_\infty$ discreto:
\[
\begin{bmatrix}
P&0&A^TP_+&C^T\\0&\eta I&B^TP_+&D^T\\
P_+A&P_+B&P_+&0\\C&D&0&I
\end{bmatrix}\succ0.
\]
Para $H_2$ contínuo, somente nas variantes $D=0$:
\[
\begin{bmatrix}-A^TP-PA-\dot P&C^T\\C&I\end{bmatrix}\succ0,
\qquad \eta\geq\operatorname{tr}(B^TPB).
\]
Para $H_2$ discreto:
\[
\begin{bmatrix}P&A^TP_+&C^T\\P_+A&P_+&0\\C&0&I\end{bmatrix}\succ0,
\qquad\eta\geq\operatorname{tr}(B^TP_+B+D^TD).
\]
No caso com parâmetros variantes, o resultado $H_2$ é interpretado como
limite uniforme da energia da resposta impulsiva, somada pelos canais de
entrada, para trajetórias admissíveis; não como norma de uma única função de
transferência LTI. Para parâmetros congelados, recupera-se a interpretação
usual de custo $H_2$. Para $H_\infty$, o certificado limita o ganho induzido
$L_2$ (ou $\ell_2$) com estado inicial nulo.

Foi subtraída margem $10^{-6}I$ de $P$ e dos blocos matriciais; as restrições
$\eta\geq0$ e de traço são não estritas. Nos ensaios Pólya, multiplicam-se
os polinômios correspondentes pela soma das coordenadas do simplex antes de
extrair os coeficientes. Esse fator vale um no domínio físico.
